
\section{ПРОЕКТУВАННЯ ТА РЕАЛІЗАЦІЯ КРОСПЛАТФОРМНОЇ СИСТЕМИ АГРЕГАЦІЇ МЕДИКО-БІОЛОГІЧНИХ ДАНИХ}

\subsection{Архітектура системи}

\subsubsection{Загальна архітектурна схема застосунку}

Розроблювана система є кросплатформним мобільним застосунком із хмарним бекендом. Система складається з двох основних частин: клієнтського мобільного застосунку, що функціонує на платформах iOS та Android, та хмарного бекенду на базі Supabase. Такий розподіл відповідає архітектурному патерну «тонкий клієнт», де уся логіка збереження та читання даних, перевірка прав доступу і бізнес-правила зосереджені на стороні сервера, а клієнт лише відображає дані та надсилає запити.

Клієнтський застосунок реалізовано засобами React Native та Expo. Він не містить власної серверної логіки та взаємодіє з хмарним бекендом виключно через \gls{REST} \gls{API} та WebSocket-підписки. Такий підхід дозволяє розгортати оновлення бекенду незалежно від оновлень застосунку в магазинах додатків, що суттєво скорочує час між виправленням серверної помилки та доступністю виправлення для кінцевих користувачів. Крім того, відсутність серверної логіки на клієнті зменшує поверхню атаки: навіть за умови компрометації клієнтського коду зловмисник не отримає доступу до даних іншого користувача.

Хмарний бекенд побудовано на платформі Supabase і містить чотири функціональних компоненти: PostgreSQL як основне сховище даних, Supabase Auth для управління автентифікацією та сесіями, Supabase Realtime для синхронізації між пристроями в реальному часі та Supabase Storage для зберігання медіафайлів. Кожен із компонентів розгортається і обслуговується платформою Supabase, що звільняє розробника від необхідності адміністрування серверної інфраструктури.

Взаємодія між клієнтом і бекендом здійснюється через \texttt{supabase-js} -- офіційний JavaScript-клієнт Supabase, що інкапсулює формування HTTP-запитів, управління \gls{JWT}-токенами та підтримку Realtime-підписок. Усі запити передаються по захищеному каналу \gls{HTTPS}; для WebSocket-з'єднань використовується протокол WSS. Завдяки типізованому клієнту з автоматично згенерованими TypeScript-інтерфейсами схеми бази даних будь-яка невідповідність між клієнтськими запитами та фактичною схемою виявляється на етапі компіляції, а не в середовищі продакшну.

При першому запуску застосунок перевіряє наявність збереженої сесії в AsyncStorage. Якщо сесія дійсна, застосунок переходить до головного екрану; в іншому випадку відображається екран автентифікації. Після успішного входу Supabase Auth видає пару токенів: \gls{JWT} доступу (термін дії 1 година) та токен оновлення. Автоматична ротація токенів відбувається при переведенні застосунку на передній план, що унеможливлює тривале використання скомпрометованого токена доступу.

\subsubsection{Архітектурні особливості клієнтської частини}

Клієнтський застосунок організовано за принципом feature-based архітектури: увесь код згруповано по функціональних доменах, а не за технічними ролями. Така структура забезпечує зв'язність коду усередині кожного модуля і мінімізує залежності між модулями, що значно спрощує навігацію по кодовій базі та локалізацію помилок. На відміну від традиційного підходу, де компоненти, хуки та утиліти одного домену розкидані по різних директоріях, feature-based підхід концентрує весь код певної функціональності в одному місці.

Директорія \texttt{app/} містить навігаційні екрани у форматі Expo Router, де структура директорій безпосередньо відображається у дерево маршрутів застосунку. Директорія \texttt{features/} зберігає доменну логіку восьми функціональних областей: feedings, sleeps, diapers, measurements, milestones, vaccinations, children та auth. Директорія \texttt{components/} містить перевикористовувані компоненти \gls{UI}, що не належать до жодного конкретного домену. Директорії \texttt{stores/}, \texttt{lib/}, \texttt{constants/} та \texttt{i18n/} зберігають глобальні стани, клієнтські інтеграції, дизайн-токени та файли локалізації відповідно. Така організація дозволяє розробнику, що не знайомий з проєктом, за назвою директорії відразу зрозуміти, яку функціональність вона реалізує, що знижує поріг входу для нових учасників команди.

Кожна директорія в \texttt{features/} містить файл \texttt{queries.ts} з React-хуками для отримання і мутації даних через TanStack Query та Supabase. Хуки отримання даних побудовані на \texttt{useQuery} і повертають стандартизовані об'єкти з полями \texttt{data}, \texttt{isLoading}, \texttt{isError} та \texttt{refetch}, що забезпечує однорідний інтерфейс взаємодії з даними по всьому застосунку. Хуки мутацій побудовані на \texttt{useMutation} і автоматично інвалідують відповідні ключі кешу після успішного виконання операції, завдяки чому \gls{UI} оновлюється без явного виклику \texttt{refetch} з боку компонентів.

Шар представлення відокремлений від шару даних: екрани у \texttt{app/} отримують дані виключно через хуки з \texttt{features/}, а не звертаються до Supabase напряму. Це дозволяє тестувати бізнес-логіку незалежно від компонентів \gls{UI} і спрощує заміну джерела даних без переписування екранів. Наприклад, якщо у майбутньому виникне необхідність замінити Supabase на інший бекенд, зміни знадобляться лише у файлах \texttt{queries.ts} кожного модуля, тоді як усі екрани і компоненти залишаться незмінними.

\subsubsection{Архітектура хмарної серверної частини}

Хмарний бекенд системи побудовано на базі платформи Supabase, що надає повний набір інфраструктурних сервісів без необхідності розробки та підтримки власного серверного застосунку. Такий підхід дозволяє зосередити весь обсяг розробки на клієнтській логіці та предметній області застосунку, не витрачаючи ресурси на розробку і підтримку типових серверних компонентів, таких як автентифікація, авторизація, файлове сховище та синхронізація в реальному часі. Усі компоненти бекенду автоматично масштабуються залежно від навантаження без участі розробника.

Центральним компонентом бекенду є PostgreSQL -- реляційна \gls{СУБД}, яка зберігає всі предметні дані системи та реалізує ізоляцію даних між користувачами через механізм \gls{RLS} (Row-Level Security), детально описаний у підрозділі~\ref{sec:rls}. Поверх PostgreSQL розгорнуто PostgREST -- проміжний шар, що автоматично генерує \gls{REST} \gls{API} для кожної таблиці і представлення бази даних. Клієнтський код звертається до бази даних через \texttt{supabase-js}, який транслює типізовані запити у відповідні HTTP-запити до PostgREST, а той виконує відповідний \gls{SQL}-запит від імені аутентифікованого користувача.

Supabase Auth реалізує управління обліковими записами і сесіями. Сервіс підтримує реєстрацію і вхід через email та пароль, а також авторизацію через Google за протоколом \gls{OAuth} 2.0. \gls{OAuth} (Open Authorization) -- відкритий стандарт делегованої авторизації, що дозволяє користувачу надати застосунку обмежений доступ до свого акаунту у стороннього провайдера (Google, Apple тощо) без передачі пароля: замість цього провайдер видає одноразовий код, який застосунок обмінює на токени доступу. При успішній автентифікації Supabase Auth генерує підписані \gls{JWT}-токени, що містять ідентифікатор користувача у полі \texttt{sub} та метадані сесії. Ці токени передаються у заголовку \texttt{Authorization: Bearer} кожного запиту до PostgREST, де PostgreSQL витягує ідентифікатор користувача функцією \texttt{auth.uid()} для застосування \gls{RLS}-фільтрів. Завдяки такій інтеграції авторизаційна перевірка відбувається безпосередньо при виконанні \gls{SQL}-запиту, а не у додатковому прошарку між клієнтом і базою даних.

Supabase Realtime транслює зміни у таблицях бази даних клієнтам через постійне WebSocket-з'єднання на основі механізму логічної реплікації PostgreSQL, забезпечуючи миттєву синхронізацію між пристроями без регулярного опитування сервера. Детальна реалізація підписок описана у підрозділі~\ref{sec:realtime}.

\subsubsection{Архітектура бази даних}

База даних системи реалізована у PostgreSQL і містить десять таблиць. Детальну \gls{ER}-діаграму бази даних наведено у додатку Д2. Всі таблиці використовують ідентифікатори типу \texttt{uuid}, що генеруються функцією \texttt{gen\_random\_uuid()}. На відміну від автоінкрементних цілочисельних ідентифікаторів, UUID не дозволяють атакуючому передбачити ідентифікатор запису іншого користувача та виконати цільований запит на його отримання.

Центральною сутністю є таблиця \texttt{children}, що зберігає профілі дітей: ідентифікатор, повне ім'я, дата народження, стать, нотатки та посилання на аватар у Supabase Storage. Зв'язок між профілями дітей та користувачами системи реалізовано через окрему таблицю \texttt{caregivers}, що містить пару (\texttt{profile\_id}, \texttt{child\_id}) та роль доглядача (\texttt{owner}, \texttt{co\_parent}, \texttt{caregiver} або \texttt{pediatrician}). Такий підхід реалізує відношення «багато-до-багатьох» між користувачами і дітьми та дозволяє одній дитині мати кількох доглядачів з різними рівнями доступу, що відповідає реальному сімейному сценарію використання застосунку.

Таблиці \texttt{feedings} та \texttt{sleeps} зберігають події з тимчасовими мітками початку та завершення; таблиця \texttt{diapers} фіксує одиничні події з єдиною міткою часу (\texttt{occurred\_at}). Таблиця \texttt{feedings} містить поле \texttt{kind} з перерахуванням допустимих типів годування (\texttt{breast\_left}, \texttt{breast\_right}, \texttt{bottle\_breast\_milk}, \texttt{bottle\_formula}, \texttt{solid}), а також поля тривалості сеансу та кількості молока в мілілітрах. Для тривалих активностей (годування та сон) реалізовано підтримку активних сеансів: запис зі значенням \texttt{NULL} у полі \texttt{ended\_at} вважається незавершеним поточним сеансом, що дозволяє відображати активний таймер на головному екрані без зберігання проміжного стану у клієнтському застосунку.

Таблиця \texttt{growth\_measurements} зберігає антропометричні вимірювання дитини з типом показника (\texttt{weight}, \texttt{height}, \texttt{head\_circumference}), числовим значенням та одиницею вимірювання. Таблиця \texttt{milestones} фіксує факти досягнення вікових міток з посиланням на шаблони (\texttt{milestone\_templates}), що зберігаються в окремій системній таблиці, заповненій стандартними педіатричними даними. Таблиця \texttt{vaccinations} зберігає факти вакцинацій без шаблонів: назву та код вакцини у текстових полях \texttt{vaccine\_name} та \texttt{vaccine\_code} безпосередньо у записі. Розподіл на таблиці шаблонів та таблиці досягнень дозволяє централізовано оновлювати набір контрольних точок розвитку без необхідності міграції індивідуальних записів кожного користувача.

\subsection{Реалізація серверної частини}

\subsubsection{Проєктування та реалізація бази даних}

Схема бази даних розроблена з урахуванням вимог нормалізації до третьої нормальної форми та особливостей механізму \gls{RLS} у Supabase. Дотримання нормалізації усуває надлишковість даних: наприклад, інформація про шаблони вікових досягнень зберігається один раз у таблиці \texttt{milestone\_templates}, а таблиця \texttt{milestones} містить лише зовнішній ключ на відповідний шаблон. Усі предметні таблиці містять первинний ключ типу \texttt{uuid} з автоматичною генерацією через \texttt{gen\_random\_uuid()}, що унеможливлює передбачення ідентифікаторів атакуючим та спрощує горизонтальне масштабування бази даних у майбутньому.

Таблиця \texttt{children} є головною сутністю схеми і містить такі поля: \texttt{id uuid PRIMARY KEY DEFAULT gen\_random\_uuid()}, \texttt{full\_name text NOT NULL}, \texttt{date\_of\_birth date NOT NULL}, \texttt{sex text CHECK (sex IN ('male', 'female', 'unspecified'))}, \texttt{notes text}, \texttt{avatar\_url text} та \texttt{created\_at timestamptz DEFAULT now()}. Обмеження \texttt{CHECK} на поле \texttt{sex} забезпечує цілісність даних на рівні \gls{СУБД} і не дозволяє клієнту зберегти довільне рядкове значення для цього поля. Таблиця \texttt{caregivers} реалізує зв'язок між користувачами та дітьми через поля \texttt{profile\_id uuid REFERENCES auth.users NOT NULL} та \texttt{child\_id uuid REFERENCES children NOT NULL} з унікальним обмеженням на цю пару, що запобігає дублюванню записів про одного й того самого доглядача.

Таблиця \texttt{feedings} містить такі ключові поля: \texttt{child\_id uuid REFERENCES children NOT NULL}, \texttt{kind text NOT NULL}, \texttt{started\_at timestamptz NOT NULL}, \texttt{ended\_at timestamptz} (зі значенням \texttt{NULL} для активного сеансу), \texttt{amount\_ml numeric(6,1)} для кількості молока у пляшковому годуванні та \texttt{created\_by uuid REFERENCES auth.users} для фіксації автора запису. Таблицю \texttt{sleeps} побудовано за аналогічним принципом; таблиця \texttt{diapers} використовує єдину мітку часу \texttt{occurred\_at} замість пари \texttt{started\_at}/\texttt{ended\_at}. У таблиці \texttt{growth\_measurements} додано поле \texttt{value numeric(7,2) NOT NULL} для числового значення та \texttt{unit text} для одиниці вимірювання, що забезпечує точність до двох десяткових знаків та підтримує зберігання як метричних, так і імперіальних одиниць.

При створенні нової дитини база даних автоматично додає запис у таблицю \texttt{caregivers} завдяки тригеру \texttt{on\_child\_created}, що виконується після кожної вставки у таблицю \texttt{children}. Тригер визначає поточного аутентифікованого користувача через \texttt{auth.uid()} та додає пару (\texttt{auth.uid()}, \texttt{new.id}) з роллю \texttt{owner}. Такий підхід гарантує, що кожна дитина завжди має хоча б одного власника, і спрощує клієнтський код: для реєстрації нової дитини потрібен лише один INSERT-запит до таблиці \texttt{children}, а зв'язок із поточним користувачем створюється автоматично на стороні \gls{СУБД} без додаткових запитів з клієнта.

\subsubsection{Реалізація Row-Level Security}
\label{sec:rls}

\gls{RLS} є ключовим механізмом ізоляції даних у системі. Після ввімкнення \gls{RLS} для таблиці PostgreSQL автоматично застосовує визначені політики до кожного запиту незалежно від того, яким способом цей запит надійшов: через \gls{API}, прямий \gls{SQL}-клієнт або будь-який інший інтерфейс. Це означає, що захист даних реалізований не на рівні логіки застосунку, а на рівні \gls{СУБД}, і не може бути обійдений навіть при помилках у клієнтському коді. Для таблиці \texttt{children} визначено дві політики. Перша дозволяє читання лише тим користувачам, для яких існує відповідний запис у таблиці \texttt{caregivers}:

\begin{verbatim}
CREATE POLICY "caregivers can view children"
  ON children FOR SELECT
  USING (
    EXISTS (
      SELECT 1 FROM caregivers
      WHERE caregivers.child_id = children.id
        AND caregivers.profile_id = auth.uid()
    )
  );
\end{verbatim}

Друга політика дозволяє вставку будь-якому аутентифікованому користувачеві, оскільки тригер \texttt{on\_child\_created} автоматично призначить його власником щойно створеного профілю дитини:

\begin{verbatim}
CREATE POLICY "authenticated users can create children"
  ON children FOR INSERT
  WITH CHECK (auth.uid() IS NOT NULL);
\end{verbatim}

Для таблиць активностей (\texttt{feedings}, \texttt{sleeps}, \texttt{diapers}, \texttt{growth\_measurements}) \gls{RLS}-політики перевіряють доступ через таблицю \texttt{caregivers} аналогічно до таблиці \texttt{children}, але умова перевіряє \texttt{child\_id} запису активності. Наприклад, для \texttt{feedings} SELECT-політика виглядає так:

\begin{verbatim}
CREATE POLICY "caregivers can view feedings"
  ON feedings FOR SELECT
  USING (
    EXISTS (
      SELECT 1 FROM caregivers
      WHERE caregivers.child_id = feedings.child_id
        AND caregivers.profile_id = auth.uid()
    )
  );
\end{verbatim}

Аналогічні політики для операцій \texttt{INSERT}, \texttt{UPDATE} та \texttt{DELETE} гарантують, що користувач може змінювати лише дані тих дітей, доглядачем яких він є. Функція \texttt{auth.uid()} повертає ідентифікатор поточного користувача з \gls{JWT}-токена, переданого в заголовку запиту; якщо токен відсутній або недійсний, функція повертає \texttt{NULL}, і жодна \gls{RLS}-політика не виконується, тобто запит від неаутентифікованого користувача автоматично повертає порожній результат для всіх захищених таблиць.

\subsection{Реалізація клієнтської частини}

\subsubsection{Організація структури модулів}

Реалізація клієнтської частини дотримується єдиного шаблону в усіх функціональних модулях: кожна директорія у \texttt{features/} реалізує повний вертикальний зріз одного домену: від типів та запитів до компонентів відображення. Шар даних кожного модуля будується на основі TanStack Query і складається з файлу \texttt{queries.ts} із хуками отримання даних (\texttt{useQuery}) та мутацій (\texttt{useMutation}), а також окремих файлів для типів і утиліт домену. Ієрархічна система ключів кешу забезпечує автоматичну інвалідацію при мутаціях без ручного виклику перезавантаження даних з боку компонентів.

Наприклад, модуль годувань містить у файлі \texttt{features/feedings/queries.ts} шість хуків: \texttt{useFeedingsForDay} та \texttt{useFeedingsInRange} для отримання записів за різними часовими діапазонами, \texttt{useActiveFeeding} для відстеження незавершеного сеансу, а також \texttt{useStartFeeding}, \texttt{useStopFeeding} та \texttt{useCreateFeeding} для операцій запису. Ключі кешу побудовано ієрархічно: \texttt{['feedings', childId]} для всіх записів дитини та \texttt{['feedings', childId, 'day', date]} для записів конкретного дня. Завдяки цій ієрархії мутація \texttt{useCreateFeeding} може одночасно інвалідувати кеш поточного дня і загальний кеш дитини одним викликом з шаблоном \texttt{['feedings', childId]}, що охоплює всі підпорядковані ключі.

Спільні компоненти \gls{UI} зосереджені в директорії \texttt{components/} і поділені на логічні групи залежно від призначення. Контейнерні компоненти \texttt{ScreenContainer} та \texttt{FormScreen} реалізують стандартні макети екранів з урахуванням безпечних зон пристрою (\texttt{SafeAreaView}) та підтримки прокручування. Компоненти форм \texttt{AppTextInput}, \texttt{DateField} та \texttt{KindGrid} забезпечують єдиний візуальний стиль полів введення даних по всьому застосунку. Графічні компоненти \texttt{DailyBarChartCard} та \texttt{MeasurementChartCard} інкапсулюють логіку побудови та відображення графіків. Усі компоненти параметризовані через TypeScript-інтерфейси та підтримують темну і світлу теми через хук \texttt{useTheme} бібліотеки React Native Paper. Для реалізації плавних анімацій у спільних компонентах використовується бібліотека \texttt{react-native-reanimated}, що виконує анімації безпосередньо у нативному потоці без передачі даних через JavaScript bridge на кожному кадрі. Завдяки цьому анімації зберігають частоту 60 кадрів на секунду навіть під час інтенсивного оновлення \gls{UI}. Reanimated застосовується у \texttt{BottomSheet} -- модальному аркуші, що виїжджає знизу екрана, де \texttt{withSpring} та \texttt{withTiming} анімують переміщення панелі по вертикалі та прозорість фону, -- та у \texttt{ActiveChildPanel}, де \texttt{withSpring} анімує індикатор активної дитини при переключенні між профілями. Спільні значення (\texttt{useSharedValue}) зберігаються у нативному потоці, а стилі обчислюються через \texttt{useAnimatedStyle} без участі JS-рантайму під час відтворення анімації.

Локалізація реалізована бібліотекою \texttt{i18next} у поєднанні з \texttt{react-i18next} із файлами перекладів у директорії \texttt{i18n/}. Застосунок підтримує дві мови: українську (за замовчуванням) та англійську. Обрана мова зберігається у AsyncStorage та застосовується при наступному запуску без необхідності повторного вибору. Всі рядки у компонентах \gls{UI} отримуються через хук \texttt{useTranslation}, що забезпечує централізоване управління текстовим контентом і спрощує додавання нових мов у майбутньому.

Статична типізація всього проєкту забезпечується TypeScript у суворому режимі (\texttt{"strict": true} у \texttt{tsconfig.json}), що вмикає перевірку на \texttt{null}/\texttt{undefined}, заборону неявних типів \texttt{any} та суворий контроль типів параметрів функцій. Псевдонім шляху \texttt{@/*} дозволяє імпортувати будь-який модуль відносно кореня проєкту без відносних шляхів виду \texttt{../../../}, що скорочує кількість помилок при переміщенні файлів. Статичний аналіз коду виконується ESLint із пресетом \texttt{eslint-config-expo}, який охоплює правила для React, React Native, TypeScript та доступності \gls{UI}-компонентів. Разом TypeScript і ESLint утворюють двошаровий бар'єр якості: компілятор відловлює структурні помилки типів, а лінтер -- проблеми стилю, небезпечні патерни та порушення правил хуків React.

\subsubsection{Моніторинг помилок та безпека середовища виконання}

Для відстеження збоїв і аномалій у продакшн-середовищі інтегровано бібліотеку \texttt{@sentry/react-native}. Sentry -- це платформа моніторингу помилок, що автоматично перехоплює необроблені винятки, збої JavaScript-рантайму та відмови нативних модулів, після чого передає структуровані звіти на хмарний дашборд із повним стек-трейсом, інформацією про пристрій, версію застосунку та контекст користувача. Ініціалізація виконується одноразово при запуску застосунку:

\begin{verbatim}
Sentry.init({
  dsn: process.env.EXPO_PUBLIC_SENTRY_DSN,
  tracesSampleRate: 0.2,
});
\end{verbatim}

Параметр \texttt{tracesSampleRate: 0.2} означає, що трасування продуктивності збирається для 20\% сесій, що дозволяє виявляти повільні мережеві запити та вузькі місця рендерингу без суттєвого навантаження на інфраструктуру. DSN-адреса (ідентифікатор проєкту Sentry) передається через змінну оточення \texttt{EXPO\_PUBLIC\_SENTRY\_DSN} і не зберігається в репозиторії.

Для захисту від запуску застосунку у скомпрометованому середовищі інтегровано бібліотеку \texttt{freerasp-react-native}. Вона виконує перевірки цілісності середовища виконання: виявляє jailbreak (iOS) та root-доступ (Android), присутність hooking-фреймворків (Frida, Xposed), запуск у емуляторі, а також модифікацію бінарного файлу застосунку. При виявленні загрози бібліотека передає подію до Sentry для фіксації у системі моніторингу та виконує налаштовану реакцію. Така інтеграція утворює дворівневий захист: Sentry фіксує помилки у штатному режимі, а freeRASP реагує на цілеспрямовані спроби маніпуляцій із застосунком.

\subsubsection{Навігаційна структура та управління станом}

Навігаційна структура застосунку реалізована засобами Expo Router, що відображає файлову систему директорії \texttt{app/} безпосередньо у дерево маршрутів. Такий підхід усуває необхідність ручної реєстрації маршрутів: додавання нового файлу \texttt{app/feedings/new.tsx} автоматично робить маршрут \texttt{/feedings/new} доступним у навігаційному стеку. Детальну структурну схему системи наведено у додатку Д3. Expo Router забезпечує також типобезпечну навігацію: при виклику \texttt{router.push('/feedings/new')} TypeScript перевіряє наявність відповідного файлу у директорії \texttt{app/}, що виключає помилки навігації на етапі компіляції.

Кореневий файл \texttt{app/\_layout.tsx} є точкою входу навігаційного стеку. Він обгортає все дерево компонентів у провайдери: \texttt{AuthProvider} для управління сесією, \texttt{QueryClientProvider} для TanStack Query та \texttt{PaperProvider} для теми матеріального дизайну. На цьому ж рівні ініціалізуються Realtime-підписки та компонент \texttt{AuthGate}, що реалізує логіку охорони маршрутів. Основний стек містить два вкладені навігатори: вкладкову навігацію з трьома головними екранами («Головна», «Статистика», «Налаштування») та модальний стек із формами введення даних, що відображаються поверх вкладок як стандартні модальні вікна платформи.

Глобальний стан застосунку зберігається у двох Zustand-сховищах з персистентністю через AsyncStorage. Сховище \texttt{activeChild} містить ідентифікатор поточно обраної дитини; при зміні цього ідентифікатора всі React Query хуки, що залежать від \texttt{childId}, автоматично виконують повторний запит до \gls{API}. Такий реактивний ланцюжок забезпечує миттєве оновлення всього інтерфейсу при переключенні між профілями дітей без явного перезавантаження будь-якого екрана. Сховище \texttt{themePreference} зберігає обраний режим теми (\texttt{system}, \texttt{light} або \texttt{dark}) і передається в \texttt{PaperProvider} для застосування відповідної теми матеріального дизайну.

\subsubsection{Зберігання файлів та синхронізація в реальному часі}
\label{sec:realtime}

Supabase Storage використовується для зберігання аватарів дітей. При завантаженні зображення застосунок обирає фото з галереї пристрою через \texttt{expo-image-picker} та стискає його до максимального розміру 512×512 пікселів для зменшення обсягу переданих даних. Стиснуте зображення завантажується у бакет \texttt{avatars} за детермінованим шляхом \texttt{\{child\_id\}/avatar.jpg}, після чого публічний \gls{URL} зображення зберігається у полі \texttt{avatar\_url} таблиці \texttt{children}. Використання детермінованого шляху замість генерованого імені файлу дозволяє автоматично перезаписувати попередній аватар при завантаженні нового без необхідності видалення старого файлу та оновлення посилання.

Синхронізація між пристроями реалізована через Supabase Realtime у файлі \texttt{lib/realtime.ts}. При активації основного стеку застосунок підписується на канал \texttt{child:\{activeChildId\}}, що відслідковує зміни у таблицях активностей конкретної дитини. Фільтр на рівні каналу обмежує трафік: клієнт отримує лише події, де \texttt{child\_id} збігається з ідентифікатором активної дитини, а не всі зміни у таблиці. При отриманні події \texttt{postgres\_changes} React Query клієнт інвалідує відповідний ключ кешу, що ініціює фонове оновлення даних:

\begin{verbatim}
supabase
  .channel(`child:${childId}`)
  .on('postgres_changes',
    { event: '*', schema: 'public', table: 'feedings',
      filter: `child_id=eq.${childId}` },
    () => queryClient.invalidateQueries(
      { queryKey: ['feedings', childId] })
  )
  .subscribe();
\end{verbatim}

При зміні активної дитини у сховищі \texttt{activeChild} попередня підписка скасовується, а нова підписка встановлюється для нового ідентифікатора. Такий підхід запобігає накопиченню неактуальних підписок та зайвому мережевому трафіку. Завдяки реалізованому механізму синхронізації, якщо один із батьків додає запис про годування зі свого пристрою, застосунок іншого батька оновлює відображення автоматично протягом 1--2 секунд без необхідності ручного оновлення сторінки.

\subsubsection{Реалізація клієнтської автентифікації}

Автентифікація реалізована через \texttt{AuthProvider} -- React-контекст, що обгортає весь застосунок і надає сесію та методи автентифікації всім дочірнім компонентам. При монтуванні провайдер викликає \texttt{supabase.auth.getSession()} для отримання збереженої сесії з AsyncStorage пристрою. Якщо сесія знайдена, провайдер перевіряє час її закінчення: якщо до закінчення залишається менше 60 хвилин, виконується автоматичне оновлення токена через \texttt{supabase.auth.refreshSession()}. У разі невдачі оновлення (наприклад, коли токен відкликано з боку сервера) сесія вважається недійсною і користувача перенаправляють на екран входу.

Для реєстрації та входу через email і пароль реалізовано відповідні мутації у файлі \texttt{features/auth/mutations.ts}. Помилки автентифікації, що повертаються Supabase Auth, перехоплюються і відображаються у формі у вигляді зрозумілих повідомлень для користувача: «Невірний email або пароль», «Email вже зареєстровано» тощо. Вхід через Google реалізовано за допомогою \gls{OAuth} 2.0 потоку з бібліотекою \texttt{expo-web-browser}: застосунок відкриває системний браузер із сторінкою авторизації Google, після успішного проходження якої браузер перенаправляє на deep-link застосунку з одноразовим кодом авторизації, який \texttt{supabase-js} автоматично обмінює на пару токенів доступу.

Компонент \texttt{AuthGate} у кореневому \texttt{\_layout.tsx} підписується на зміни стану автентифікації через \texttt{supabase.auth.onAuthStateChange} та умовно відображає або основний стек із вкладками, або стек автентифікації. Завдяки реактивній підписці стан навігації оновлюється миттєво при зміні сесії: як при вході, так і при виході. Екран автентифікації побудований за принципом єдиної точки входу: один екран обслуговує як вхід, так і перехід до реєстрації. Верхня частина екрана є брендовою зоною з градієнтним фіолетовим фоном та логотипом застосунку, нижня містить білу картку з двома варіантами автентифікації: формою email/пароль та кнопкою «Продовжити з Google». Перемикач мови у правому верхньому куті дозволяє змінювати локалізацію безпосередньо на екрані входу без необхідності авторизації. Після успішної реєстрації застосунок автоматично виконує вхід і перенаправляє на екран створення першого профілю дитини.

\begin{verbatim}
supabase.auth.getSession().then(({ data }) => {
  setSession(data.session);
  setLoading(false);
});
const { data: sub } = supabase.auth.onAuthStateChange(
  (_event, s) => setSession(s)
);
return () => sub.subscription.unsubscribe();
\end{verbatim}

\subsubsection{Реалізація ключових екранів та компонентів}

Головний екран застосунку (\texttt{app/(tabs)/index.tsx}) є центральним інформаційним хабом і першим, що бачить користувач після входу. У верхній частині відображається компонент \texttt{HeroCard} -- вітальна картка з іменем поточного користувача та датою; нижче -- картка активної дитини з аватаром, ім'ям та кількістю повних місяців від дати народження. Рядок лічильників \texttt{StatsRow} відображає три агреговані показники поточного дня: кількість годувань, загальний час сну та кількість змін підгузків. Якщо є активний незавершений сеанс, відображається також компонент \texttt{ActiveTimerCard} з таймером у реальному часі та кнопкою зупинки. Сітка швидких дій надає прямий доступ до форм введення всіх типів активностей, а хронологічна стрічка подій у нижній частині агрегує записи з усіх доменів у зворотному хронологічному порядку. Навігація по дням реалізована через кнопки «назад» та «вперед», що змінюють значення \texttt{selectedDay} через \texttt{subDays}/\texttt{addDays} з \texttt{date-fns}, дозволяючи переглядати стрічку будь-якого минулого дня без переходу на окремий екран.

Екран статистики (\texttt{app/(tabs)/stats.tsx}) відображає аналітику за останні 7 днів у вигляді трьох стовпчастих графіків: кількість годувань, загальний час сну та кількість змін підгузків по днях тижня. Компонент \texttt{DailyBarChartCard} для кожного показника отримує масив значень за 7 днів і будує графік засобами бібліотеки \texttt{react-native-gifted-charts}, де кожен стовпець відповідає одному дню. Для аналізу режиму сну реалізовано окремий компонент \texttt{SleepPatternChart}, де вісь Y відображає години доби від 00:00 до 24:00, а стовпці показують діапазон активного сну у кожен день, дозволяючи візуально виявляти порушення нічного режиму. Під графіком сну відображається підсумковий рядок із тривалістю денного та нічного сну за тиждень. Якщо середня кількість замін підгузків виходить за межі норми (менше 5 або більше 8 на добу), екран відображає відповідне попередження лікарського характеру.

Екран статистики підтримує експорт даних у формат PDF через функцію \texttt{exportStatsPdf} у файлі \texttt{lib/pdf.ts}. Функція динамічно завантажує бібліотеки \texttt{expo-print} та \texttt{expo-sharing} через \texttt{require()} замість статичного імпорту, що дозволяє екрану статистики коректно завантажуватись у Expo Go, де ці нативні модулі відсутні. При активації експорту функція \texttt{buildHtml} генерує повноцінний HTML-документ із CSS-стилями, який містить вітальну секцію з іменем дитини та датою, окремі картки для годувань, сну та підгузків із таблицями по днях, а також секцію антропометричних вимірювань. Бібліотека \texttt{expo-print} конвертує цей HTML у \texttt{.pdf}-файл засобами нативного рушія платформи (\texttt{Print.printToFileAsync}), а \texttt{expo-sharing} відкриває стандартний системний діалог передачі файлу (\texttt{Sharing.shareAsync}), через який користувач може зберегти PDF у Files, надіслати електронною поштою або передати через AirDrop. Всі текстові рядки та формати дат у PDF локалізовані відповідно до поточної мови застосунку через \texttt{i18next} та \texttt{date-fns}.

Форми введення даних побудовані на єдиному компоненті \texttt{FormScreen}, що реалізує стандартний макет із заголовком, кнопкою закриття та прокручуваним тілом з урахуванням клавіатури. Вибір типу активності реалізований через компонент \texttt{KindGrid} -- сітку іконок-плиток, де кожен варіант відображається у вигляді картки з кольоровою іконкою та підписом, а активний елемент підсвічується кольоровою рамкою. Форма годування підтримує п'ять типів: ліву та праву груди, пляшку з грудним молоком, пляшку з сумішшю і тверду їжу, -- а також два режими введення: запуск таймера в реальному часі кнопкою «Старт зараз» або ручне введення тривалості минулого сеансу. Форма підгузків підтримує чотири типи: мокрий (\texttt{wet}), брудний (\texttt{dirty}), змішаний (\texttt{mixed}) та сухий (\texttt{dry}), -- що дозволяє точніше відстежувати стан дитини та виявляти можливі проблеми зі здоров'ям. Вибір дати і часу реалізований через \texttt{DateField}, що відкриває нативний системний пікер дати платформи, забезпечуючи звичний для користувача досвід.

\begin{verbatim}
export function useStartFeeding() {
  const qc = useQueryClient();
  return useMutation({
    mutationFn: async ({ childId, kind }) => {
      const { data, error } = await supabase
        .from('feedings')
        .insert({ child_id: childId, kind,
                  started_at: new Date().toISOString(),
                  ended_at: null })
        .select('*').single();
      if (error) throw error;
      return data;
    },
    onSuccess: (f) =>
      qc.invalidateQueries({ queryKey: feedingsKey(f.child_id) }),
  });
}
\end{verbatim}

Екран налаштувань (\texttt{app/(tabs)/settings.tsx}) організований у три секції. Секція «Профіль» надає доступ до редагування імені та email, а також до зміни пароля через відповідні модальні форми. Секція «Налаштування» містить перемикач мови (українська/англійська) та вибір теми оформлення. Секція «Акаунт» включає кнопки виходу та видалення облікового запису, кожна з яких відкриває \texttt{ConfirmDialog} для підтвердження деструктивної операції. Зміна мови відбувається реактивно через хук \texttt{useLocale} без перезавантаження застосунку.

Екран вікових досягнень (\texttt{app/milestones/index.tsx}) відображає шаблони контрольних точок розвитку дитини, згрупованих за віковими діапазонами. Шаблони завантажуються з таблиці \texttt{milestone\_templates} і збагачуються даними про фактичні досягнення дитини з таблиці \texttt{milestones}. Якщо вік дитини потрапляє у діапазон (\texttt{age\_min\_months}--\texttt{age\_max\_months}) шаблону, і досягнення ще не зафіксоване, шаблон потрапляє до секції «Актуально зараз» на початку екрана, привертаючи увагу батьків до поточних цілей розвитку. Кожен шаблон відображається як картка з категорійною іконкою, описом та статусом; натискання відкриває екран деталей, де можна зафіксувати дату досягнення та додати нотатку.

Екран вакцинацій (\texttt{app/vaccinations/index.tsx}) побудований аналогічно до екрана досягнень, але групує щеплення за педіатричними слотами (при народженні, 2 місяці, 4 місяці, 6 місяців тощо). Записи завантажуються з таблиці \texttt{vaccinations}, де зберігаються код та назва вакцини, дата введення (\texttt{administered\_at}) та дата планового введення (\texttt{scheduled\_for}). Секція «Актуально зараз» виокремлює записи, термін яких настав відповідно до поточного віку дитини. Кожна картка відображає назву вакцини та статус введення; при натисканні відкривається форма фіксації дати введення та нотатки.

Редагування профілю дитини доступне через екран \texttt{app/children/[id]/edit.tsx}, що дозволяє змінити ім'я, дату народження, стать та нотатки, а також завантажити або видалити аватар через \texttt{expo-image-picker}. Видалення акаунту виконується з екрана налаштувань і вимагає підтвердження через \texttt{ConfirmDialog}, після чого каскадно видаляються всі пов'язані записи завдяки \texttt{ON DELETE CASCADE} у схемі бази даних.

\subsubsection{Робота з датами та часом}

Уся робота з датами та часом у застосунку виконується бібліотекою \texttt{date-fns} -- модульною утилітою, що надає функції для форматування, арифметики та порівняння дат без мутації об'єктів і без глобального стану. На відміну від \texttt{moment.js}, \texttt{date-fns} підтримує tree-shaking: до фінального пакету потрапляють лише ті функції, що реально імпортуються, що зменшує розмір збірки. Бібліотека використовується у трьох ключових контекстах. По-перше, у функції \texttt{formatAge} (\texttt{lib/age.ts}) для обчислення віку дитини: \texttt{differenceInDays}, \texttt{differenceInMonths} та \texttt{differenceInYears} визначають точний вік і обирають правильну одиницю відображення (дні, тижні, місяці або роки). По-друге, у функції \texttt{aggregateByDay} (\texttt{lib/aggregate.ts}) для формування 7-денних бакетів статистики: \texttt{subDays} генерує масив дат від сьогодні назад, \texttt{parseISO} перетворює рядки з бази даних у об'єкти \texttt{Date}, а \texttt{isSameDay} зіставляє записи активностей із відповідним бакетом. По-третє, у модулі PDF-експорту (\texttt{lib/pdf.ts}) для локалізованого форматування дат у звіті через \texttt{format} із передачею об'єкта \texttt{locale} відповідно до поточної мови інтерфейсу.

\subsubsection{Реалізація алгоритмів обробки медичних даних}

Алгоритм класифікації антропометричних показників відповідно до стандартів \gls{ВООЗ} реалізовано у файлі \texttt{features/measurements/who/index.ts}. Вихідні перцентильні таблиці зберігаються статично у \texttt{who/standards.ts} у вигляді вкладених числових масивів формату \texttt{[age, P3, P15, P50, P85, P97]}, де \texttt{age} -- вік у місяцях окремо для кожної статі та кожного типу показника. Такий підхід дозволяє уникнути мережевих запитів при кожному розрахунку, оскільки таблиці компілюються безпосередньо у пакет застосунку та залишаються незмінними між сесіями. Стандарти \gls{ВООЗ} охоплюють вікові діапазони від 0 до 24 місяців для дівчаток та хлопчиків окремо, що дозволяє враховувати статеві відмінності у нормах фізичного розвитку.

Функція \texttt{whoBands(sex, kind, ageMonths)} виконує лінійну інтерполяцію між двома суміжними рядками таблиці для точного вікового значення дитини. Спочатку обирається відповідна таблиця за статтю та типом показника, потім знаходяться два сусідні рядки, між якими знаходиться вік дитини (\texttt{floor(ageMonths)} та \texttt{ceil(ageMonths)}). Для кожного з п'яти перцентилів (P3, P15, P50, P85, P97) обчислюється інтерпольоване значення за формулою:
\[ v = v_0 + (v_1 - v_0) \times (ageMonths - month_0), \]
де $v_0$ та $v_1$ -- значення перцентиля для нижнього та верхнього рядків таблиці, $month_0$ -- нижній цілий вік у місяцях. Функція повертає об'єкт \texttt{BandValues} з інтерпольованими значеннями п'яти перцентилів, що використовуються для побудови зон на графіку та класифікації вимірювання.

Функція \texttt{classifyValue(value, bands)} визначає, в яку зону потрапляє виміряне значення: нижче P3, між P3 і P15, між P15 і P85 (норма), між P85 і P97, або вище P97. Результати класифікації відображаються на графіку кольоровими зонами: зелена відповідає нормі між P15 та P85, жовті -- граничним значенням P3--P15 і P85--P97, червоні -- патологічним відхиленням нижче P3 та вище P97. Текстовий підпис під графіком дублює класифікацію словесно, що полегшує сприйняття для батьків без медичної освіти.

\begin{verbatim}
export function classifyValue(
  value: number, bands: Bands
): PercentileBand {
  if (value < bands.p3)  return 'below_p3';
  if (value < bands.p15) return 'p3_p15';
  if (value <= bands.p85) return 'p15_p85';
  if (value <= bands.p97) return 'p85_p97';
  return 'above_p97';
}
\end{verbatim}

Алгоритм агрегації даних сну реалізовано у хуку \texttt{useSleepPattern}. Хук отримує список сесій сну за 7 днів і для кожного дня обчислює два показники: тривалість денного та нічного сну. Нічним вважається сон, що частково або повністю знаходиться в проміжку 19:00--06:59 наступного дня. Тривалість частини кожного сну, що перетинається з нічним вікном, обчислюється як довжина перетину двох часових відрізків, після чого значення підсумовуються для кожного дня окремо. Отримані дані передаються до компонента \texttt{SleepPatternChart}, який будує двовимірне подання режиму сну, що дозволяє батькам оцінити регулярність та тривалість нічного відпочинку дитини протягом тижня і своєчасно виявити порушення режиму сну.

\begin{conclusionsection}

\section{ВИСНОВОК ДО РОЗДІЛУ 3}

У третьому розділі розглянуто реалізацію кросплатформної клієнт-серверної системи агрегації та обробки медико-біологічних даних.

Спроектовано та реалізовано архітектуру системи, що включає клієнтський застосунок на базі React Native та Expo з feature-based організацією коду і хмарний бекенд на платформі Supabase. Відокремлення шару представлення від шару даних та зосередження логіки авторизації на рівні \gls{СУБД} забезпечують підтримуваність коду та унеможливлюють витік даних між обліковими записами незалежно від логіки клієнтського застосунку.

Реалізовано реляційну схему бази даних PostgreSQL з десятьма таблицями відповідно до вимог третьої нормальної форми. Для кожної таблиці визначено \gls{RLS}-політики, що ізолюють дані різних користувачів на рівні \gls{СУБД} без додаткової логіки авторизації в клієнтському коді. Реалізовано тригер автоматичного призначення власника при створенні нового профілю дитини, що скорочує кількість клієнтських запитів із двох до одного.

Реалізовано систему автентифікації з підтримкою email/пароль та Google OAuth, а також синхронізацію між пристроями в реальному часі через Supabase Realtime із середньою затримкою 1--2 секунди.

На клієнтській частині реалізовано алгоритм класифікації антропометричних показників відповідно до стандартів \gls{ВООЗ} з лінійною інтерполяцією перцентильних таблиць для точного вікового значення дитини, а також алгоритм аналізу структури сну з розподілом на денну та нічну складові. Розроблені алгоритми медично коректні та відповідають офіційним рекомендаціям \gls{ВООЗ} щодо оцінювання фізичного розвитку дітей.

\end{conclusionsection}
